\section{分类、神经网络与信息论边界}
\label{sec:07-Information-Theory/06-Information-and-Learning/04}

训练日志停在验证交叉熵 \(0.62\) 上不动了。换更大的网络，多训五十轮，加数据增强——曲线纹丝不动。此刻团队里会出现两种声音：一种说模型还不够强，另一种说特征里根本就没有更多东西可挖。

这两种诊断指向完全相反的行动。前者要加算力，后者要加数据源。猜错的代价是一个季度。信息论对这件事有话说，而且说得比大多数经验法则硬。

\subsection{交叉熵损失的两栏账}

分类模型输出的是条件分布 \(q_\theta(y\given x)\)，总体上的平均负对数损失是

\[
\mathcal R(\theta)=\E_{(X,Y)\sim P}\bigl[-\log q_\theta(Y\given X)\bigr].
\]

在期望号里加减 \(\log P(Y\given X)\)，把被积项拆成 \(-\log P(Y\given X)\) 和对数比两块，分别取期望，就得到

\[
\mathcal R(\theta)=H_P(Y\given X)+\E_X\bigl[\KL(P_{Y\given X}\,\|\,q_\theta(\cdot\given X))\bigr].
\]

这是\hyperref[sec:07-Information-Theory/03-Divergences/01]{``错误模型的代价：从交叉熵到 KL 散度''}那个分解的条件版本，读法完全一样。第一项是看过 \(X\) 之后标签仍然剩下的不确定度，跟 \(\theta\) 没有关系；第二项是模型条件分布与真实条件分布的失配。训练能动的只有第二项，第一项是地板。

用上一节和上上节的语言说，这就是一份码长账单：总的平均码长里，有一栏是数据自身的，另一栏是用错分布多付的。梯度下降在削的一直是后一栏。

\begin{center}
\begin{tikzpicture}[x=1.25cm,y=0.82cm,font=\small]
  \fill[blue!16] (0,0) rectangle (6.15,0.9);
  \fill[red!13] plot[domain=0:6,samples=70] (\x,{0.9+2.6*exp(-\x/1.4)})
    -- (6,0.9) -- (0,0.9) -- cycle;
  \draw[->,black!55] (0,0) -- (6.7,0) node[right,font=\footnotesize] {训练进程};
  \draw[->,black!55] (0,0) -- (0,4.1) node[above,font=\footnotesize] {平均交叉熵};
  \draw[black!45,dashed] (0,0.9) -- (6.5,0.9);
  \draw[blue!80,very thick] plot[domain=0:6,samples=70] (\x,{0.9+2.6*exp(-\x/1.4)});
  \draw[black!55,dashed,thick] plot[domain=0:6,samples=70]
    (\x,{0.9+2.6*exp(-\x/1.2)+0.045*\x*\x});
  \node[right,font=\footnotesize,black!75] at (4.35,1.28) {训练};
  \node[right,font=\footnotesize,black!75] at (5.15,2.65) {验证};
  \node[right,font=\footnotesize,black!70] at (2.15,1.95) {\(\E_X\KL\)：可优化};
  \node[right,font=\footnotesize,black!70] at (2.15,0.42) {\(H(Y\given X)\)：地板};
\end{tikzpicture}

\emph{红色区域是模型自己欠下的失配，训练削的就是它。蓝色地板由数据决定，加多少参数都不会下沉。虚线提醒另一件事：训练损失贴着地板并不表示验证损失也在下降，两条曲线分开的地方是过拟合，与地板高低无关。}
\end{center}

单个样本上，one-hot 标签让损失塌缩成 \(-\log q_\theta(y\given x)\)，这解释了为什么``自信地错''会被狠狠罚——真实类别的预测概率从 \(0.9\) 掉到 \(0.1\)，损失从约 \(0.105\) 涨到约 \(2.303\)。

顺带把困惑度接上。交叉熵是平均码长，报出来是对数尺度，指数回去就得到等效候选数。语言模型的损失以 nat 为单位，取 \(e\) 的指数得到的困惑度读作``模型平均每一步在多少个词之间犹豫''。困惑度从 \(50\) 降到 \(20\)，就是等效候选从五十个缩到二十个。它有一条不会被突破的底线：模型算的是交叉熵，交叉熵不小于条件熵，所以困惑度不可能降到 \(H(Y\given X)\) 对应的那个值以下。换算的细节在\hyperref[sec:07-Information-Theory/01-Information-and-Entropy/02]{``Shannon 熵与二元熵''}。

\subsection{Fano：条件熵给错误率设的硬下界}

地板的存在解释了损失为什么压不下去，但还没有换算成大家真正关心的指标。Fano 不等式做的就是这个换算。

\begin{theorem}[Fano 不等式]
设类别 \(Y\in\set{1,\dots,K}\)，预测 \(\hat Y=g(X)\) 是特征的函数，错误率 \(P_e=P(\hat Y\ne Y)\)。则
\[
H(Y\given X)\le h_2(P_e)+P_e\log_2(K-1).
\]
\end{theorem}

条件对账值得做一遍。这里要求 \(\hat Y\) 只经由 \(X\) 产生，也就是 \(Y\to X\to\hat Y\) 构成 Markov 链——若预测过程偷看了标签，不等式当然失效，而这正是信息泄漏的表现形式。\(K\) 必须有限，右端的 \(\log_2(K-1)\) 才有意义。除此之外不需要任何分布形状的假设，也不管 \(g\) 是线性模型还是千亿参数的网络。

论证的骨架是一次链式法则加一次放缩，不必展开。引入错误指示 \(E=\ind\set{Y\ne\hat Y}\)：知道 \(\hat Y\) 之后，``有没有错''最多值 \(h_2(P_e)\) 比特；一旦错了，还要在剩下的 \(K-1\) 类里定位，这部分按错误发生的概率折算，最多值 \(P_e\log_2(K-1)\) 比特。两项相加已经覆盖了 \(Y\) 在给定 \(\hat Y\) 之后的全部不确定度，而 \(H(Y\given X)\le H(Y\given\hat Y)\)，因为 \(\hat Y\) 是 \(X\) 的函数，条件熵只会因为条件变多而减小。

把不等式反过来解，粗略地用 \(h_2(P_e)\le 1\) 并把 \(\log_2(K-1)\) 放大成 \(\log_2 K\)，得到一个更好用的形式：

\[
P_e\ge\frac{H(Y\given X)-1}{\log_2 K}.
\]

\begin{center}
\begin{tikzpicture}[x=1.9cm,y=2.35cm,font=\small]
  \fill[red!13] (1,0) -- (3.322,0.699) -- (3.322,0) -- cycle;
  \draw[->,black!55] (0,0) -- (3.75,0)
    node[right,font=\footnotesize] {\(H(Y\given X)\)（比特）};
  \draw[->,black!55] (0,0) -- (0,1.12) node[above,font=\footnotesize] {\(P_e\)};
  \draw[red!75,very thick] (1,0) -- (3.322,0.699);
  \draw[black!30,dashed] (3.322,0) -- (3.322,1.02);
  \node[above,font=\footnotesize,black!60] at (3.322,1.0) {\(\log_2 K\)};
  \draw[black!35,dashed] (2.0,0) -- (2.0,0.301) -- (0,0.301);
  \fill[black!70] (2.0,0.301) circle (1.3pt);
  \node[left,font=\footnotesize] at (-0.04,0.301) {\(0.30\)};
  \node[left,font=\footnotesize] at (-0.04,1.0) {\(1\)};
  \node[below,font=\footnotesize] at (0,-0.04) {\(0\)};
  \node[below,font=\footnotesize] at (2.0,-0.04) {\(2\)};
  \node[font=\footnotesize,black!65,align=center] at (2.55,0.22) {没有任何\\分类器能进来};
  \node[font=\footnotesize,black!60,align=center] at (1.55,0.78)
    {界不排除这一片，\\但也不保证够得着};
  \node[font=\footnotesize,black!55,align=center] at (0.5,0.66) {界退化为\\\(P_e\ge0\)};
\end{tikzpicture}

\emph{取 \(K=10\)。特征看完之后若标签还剩 \(2\) 比特的条件熵，任何算法的错误率都不会低于约 \(30\%\)——红色三角是禁区，与模型架构无关。左侧那一段是这个放松版本失效的区域：条件熵不足 \(1\) 比特时右端为负，界只剩下平凡的 \(P_e\ge0\)。}
\end{center}

要注意这是下界，不是预测。它说的是``不可能比这更好''，没有说``一定能做到这么好''；真正的 Bayes 错误率取决于后验分布的具体形状，而 Fano 只用了一个标量，所以通常不紧。粗放缩之后的版本尤其松，右端为负时它什么也没说。

\subsection{分不清是模型不够好还是信息不够，怎么办}

理论上的判据很干净：把 \(H(Y\given X)\) 算出来，代进 Fano。实践上这条路基本走不通——高维特征下直接估计互信息或条件熵极不可靠，估计器之间的差异常常比想测的效应还大。所以不要指望把 Fano 界算成一个验收指标。它的价值在于提示排查方向，而排查有几种更实在的做法。

最直接的是量标签噪声。让多个标注者在同样的特征下独立打标，统计他们彼此的不一致率。如果专家之间就有 \(15\%\) 的分歧，那说明在这套特征下标签本身带着可观的条件熵，指望模型做到 \(95\%\) 准确没有依据。这个不一致率不等于 Bayes 错误率，但它是一个能实测、能复现的参照，而且往往比模型调优的空间大得多。

其次是画两条学习曲线。横轴放训练集比例，从 \(10\%\) 走到 \(100\%\)：测试误差若还在稳定下降，缺的是样本；若已经贴着一条水平线，加数据帮不上忙。横轴换成模型容量再走一遍：误差随容量增大而下降说明模型受限，早早压平则说明瓶颈在别处。两条曲线都压平在同一个正的水平上，这才是``信息不够''的典型证据。这套诊断与\hyperref[sec:13-Machine-Learning/10-Learning-Theory/01]{``训练误差为何总是乐观''}讨论的乐观偏差是一体两面。

第三是用非参数方法探底。样本充足时，\(k\) 近邻的测试误差可以当作 Bayes 错误率的粗糙探针——最近邻分类器的渐近错误率不超过 Bayes 错误率的两倍，这给了一个能实测的上界包围。它在高维会退化，所以只能当线索，不能当结论。

还有一条最省事的检查：换一批特征。如果加进一个新的数据源之后地板明显下沉，那就证明原来的地板不是任务固有的，而是特征选择造成的。地板是相对于给定特征而言的条件熵，不是关于任务的常数——这一点比 Fano 不等式本身更常被误读。

\subsection{堆层数不会造出标签信息}

对参数已固定、每层表示只由上一层生成的前馈网络 \(X\to T_1\to\cdots\to T_L\)，数据处理不等式逐层生效：

\[
I(Y;T_{k+1})\le I(Y;T_k)\le I(Y;X).
\]

这看着和``深层特征越来越好用''打架，其实比较的不是同一件事。后面的层不能增加总体上可用的标签信息，却可以把信息重排成线性分类器读得懂的形状。一个可逆旋转不改变互信息，却完全改变了坐标轴的对齐方式；一个有损的非线性可能丢掉少量无关细节，同时把决策边界削得简单许多。网络变强不需要互信息逐层上升——信息量、可提取性、几何结构和解码器复杂度是四个不同的问题。完整讨论在\hyperref[sec:07-Information-Theory/02-Joint-Information/04]{``数据处理不等式与充分表示''}。

同样的道理反过来用，可以识破一类特征工程的自我欺骗。若经验估计报出

\[
I(Y;T)>I(Y;X),
\]

先查的应该是估计偏差、采样误差和变量定义，而不是宣布神经网络推翻了数据处理不等式。真出现这种读数，第二个要查的是标签泄漏——特征的生成过程有没有直接接触过目标变量。

单变量排序在这里也会漏掉模式，异或就是最小反例：\(X_1\) 与 \(X_2\) 各自与标签的互信息都是零，联合却完全决定标签。只保留单变量信息增益最大的特征，可能一个都选不中，而非线性网络能把两者组合出有效表示。另一头是重复特征：两个高度相关的列各自互信息都很高，第二个带来的条件互信息却近乎为零。独特信息、冗余信息与协同信息是三回事，一个标量排序描述不了多变量结构。

\subsection{几个常被当成同义词的判断}

最后把容易混用的几个说法分开。交叉熵低意味着模型给真实标签分配了较高概率，它同时受概率质量和置信度影响；准确率高只说最大概率的那一类经常正确，剩下的概率怎么分配它不管；校准好则是另一件事——预测为 \(70\%\) 的那批事件长期确实约有 \(70\%\) 发生，这一项与前两项都可以脱钩，评价方式见\hyperref[sec:11-Mathematical-Statistics/07-Statistical-Decision-and-Reliable-Prediction/04]{``Proper scoring rule 与校准衡量概率是否诚实''}。至于 \(I(Y;T)\) 高和 \(I(X;T)\) 低，它们是在总体分布和一组明确的随机变量定义下才成立的陈述，落到实测就要面对上一节说的那些估计陷阱。

这几项彼此相关，但没有一项能替另一项作证。实验里该分别测量，而不是拿``信息更多''一句话概括。

本单元到此把信息论能对学习说的话摆完了，它们的形态是一致的：都是关于上限和下限的陈述，都建立在``随机变量是什么、概率模型是什么''被交代清楚之后。它们不承诺任何算法能达到边界，也不区分关联与因果——这两件事得靠别的工具。想看同一批概念在机器学习自己的语境里怎样展开，可以接着读\hyperref[ch:119]{``信息论学习：把预测、约束、复杂度与表示放到同一尺度''}。

下一部分换一个方向的可信性问题。这里假定了公式一旦写对，数值就跟着对；而浮点数、条件数与迭代格式会告诉你，写对的公式在机器上照样可能算错。
